% !TEX program = xelatex
\documentclass[11pt,a4paper]{article}

\usepackage[UTF8,fontset=none]{ctex}
\setCJKmainfont[BoldFont=Heiti SC,ItalicFont=Songti SC]{Songti SC}
\setCJKsansfont{Heiti SC}
\setCJKmonofont{Heiti SC}

\usepackage{amsmath,amssymb,amsfonts,mathtools}
\usepackage{booktabs,array,tabularx}
\usepackage[margin=1.9cm]{geometry}
\usepackage[dvipsnames]{xcolor}
\usepackage[most]{tcolorbox}
\usepackage[hidelinks]{hyperref}
\usepackage{enumitem}
\usepackage{fancyhdr}

\setlength{\parindent}{0em}
\setlength{\parskip}{0.45em}
\setlist[itemize]{leftmargin=*,topsep=2pt,itemsep=1.5pt}
\setlist[enumerate]{leftmargin=*,topsep=2pt,itemsep=1.5pt}
\allowdisplaybreaks

\pagestyle{fancy}
\fancyhf{}
\lhead{\small Statistical Modelling and Inference}
\rhead{\small 2026 模拟样卷与详细解析}
\cfoot{\thepage}

\hypersetup{
  pdftitle={Statistical Modelling and Inference 2026 Sample Examination with Solutions},
  pdfauthor={Prepared from SMI Final Review 2026}
}

\newcommand{\E}{\operatorname{E}}
\newcommand{\Var}{\operatorname{Var}}
\newcommand{\Cov}{\operatorname{Cov}}
\newcommand{\Bias}{\operatorname{Bias}}
\newcommand{\MSE}{\operatorname{MSE}}
\newcommand{\SE}{\operatorname{SE}}
\newcommand{\iid}{\overset{\mathrm{iid}}{\sim}}
\newcommand{\by}{\mathbf{y}}
\newcommand{\bX}{\mathbf{X}}
\newcommand{\bH}{\mathbf{H}}
\newcommand{\bI}{\mathbf{I}}
\newcommand{\be}{\mathbf{e}}
\newcommand{\bbeta}{\boldsymbol{\beta}}
\newcommand{\bvarepsilon}{\boldsymbol{\varepsilon}}
\newcommand{\pts}[1]{\hfill\textbf{[#1 分]}}

\newtcolorbox{instructionbox}[1][]{
  colback=gray!5!white,colframe=gray!65!black,
  fonttitle=\bfseries,title=#1,breakable
}
\newtcolorbox{questionbox}[1][]{
  colback=blue!3!white,colframe=blue!60!black,
  fonttitle=\bfseries,title=#1,breakable
}
\newtcolorbox{solutionbox}[1][]{
  colback=green!3!white,colframe=green!55!black,
  fonttitle=\bfseries,title=#1,breakable
}
\newtcolorbox{checkbox}[1][]{
  colback=orange!5!white,colframe=orange!70!black,
  fonttitle=\bfseries,title=#1,breakable
}
\newtcolorbox{warnbox}[1][]{
  colback=red!4!white,colframe=red!60!black,
  fonttitle=\bfseries,title=#1,breakable
}

\title{\textbf{Statistical Modelling and Inference}\\[0.3em]
2026 期末模拟样卷与详细解析}
\author{依据确认后的期末范围与复习讲义编制}
\date{June 2026}

\begin{document}
\maketitle

\begin{instructionbox}[考试说明]
\begin{itemize}
  \item 建议时长：120 分钟；满分：100 分；共五道大题，每题可拆成若干小问。
  \item 本卷覆盖：无偏估计与 MSE、置信区间与临界值法检验、简单线性回归手算、Week 9 回归证明、残差性质、四类分布的 MLE 推断。
  \item 本卷不涉及：ANOVA、R 编程、P-value、Power、polynomial regression、nested-model comparison。
  \item 除非另有说明，显著性水平均为 $\alpha=0.05$，保留至少三位小数。
\end{itemize}
\end{instructionbox}

\begin{instructionbox}[本卷可直接使用的临界值]
\[
  z_{0.05}=1.645,
  \qquad
  z_{0.025}=1.960.
\]
\[
\begin{aligned}
  t_{0.05,15}&=1.753,
  &t_{0.025,15}&=2.131,\\
  t_{0.025,20}&=2.086,
  &t_{0.025,6}&=2.447.
\end{aligned}
\]
\end{instructionbox}

\tableofcontents
\newpage

% ============================================================
\section{模拟试卷}
% ============================================================

\begin{questionbox}[第 1 题：无偏估计与 MSE（14 分）]
设 $T_1,T_2$ 是参数 $\theta$ 的两个相互独立无偏估计量，并且
\[
  \Var(T_1)=\frac4n,
  \qquad
  \Var(T_2)=\frac9n.
\]
考虑线性组合
\[
  T_a=aT_1+(1-a)T_2,
\]
其中 $a$ 为实数。

\begin{enumerate}
  \item 证明一般恒等式
  \[
    \MSE(T)=\Var(T)+\Bias(T)^2.
  \]
  \pts{4}
  \item 证明 $T_a$ 对任意 $a$ 都是 $\theta$ 的无偏估计量，并求使 $\MSE(T_a)$ 最小的 $a$。 \pts{6}
  \item 求最小 MSE，并与单独使用 $T_1$ 或 $T_2$ 比较。解释最优权重为何更偏向 $T_1$。 \pts{4}
\end{enumerate}
\end{questionbox}

\begin{questionbox}[第 2 题：单总体均值、未知方差与检验错误（16 分）]
某总体服从正态分布，其均值 $\mu$ 与方差均未知。随机抽取 $n=16$ 个观测，得到
\[
  \bar x=52.0,
  \qquad
  s=4.0.
\]

\begin{enumerate}
  \item 构造 $\mu$ 的 95\% 置信区间。 \pts{6}
  \item 使用临界值法检验
  \[
    H_0:\mu\leq50
    \qquad\text{vs.}\qquad
    H_a:\mu>50.
  \]
  写出统计量、参考分布、拒绝域、决定和语境结论。 \pts{7}
  \item 在本题语境下分别描述第一类错误和第二类错误。 \pts{3}
\end{enumerate}
\end{questionbox}

\begin{questionbox}[第 3 题：两总体均值与总体比例（18 分）]
\textbf{第一部分：两个独立均值。}

从两个相互独立的正态总体中抽样，并可假设总体方差相等。样本摘要为
\[
\begin{array}{c|ccc}
 & n & \bar x & s\\
\hline
\text{总体 1} & 10 & 78 & 6\\
\text{总体 2} & 12 & 72 & 5
\end{array}
\]

\begin{enumerate}
  \item 计算 pooled variance、pooled standard deviation 以及
  $\SE(\bar X_1-\bar X_2)$。 \pts{4}
  \item 构造 $\mu_1-\mu_2$ 的 95\% 置信区间。 \pts{3}
  \item 使用临界值法检验
  \[
    H_0:\mu_1-\mu_2=0
    \qquad\text{vs.}\qquad
    H_a:\mu_1-\mu_2\ne0.
  \]
  \pts{3}
\end{enumerate}

\textbf{第二部分：单总体比例。}

随机抽取 200 人，其中 116 人支持某项提议。令 $p$ 表示总体支持比例。
\begin{enumerate}[resume]
  \item 构造 $p$ 的 95\% Wald 置信区间。 \pts{4}
  \item 使用临界值法检验
  \[
    H_0:p\leq0.5
    \qquad\text{vs.}\qquad
    H_a:p>0.5.
  \]
  注意检验标准误与置信区间标准误的区别。 \pts{4}
\end{enumerate}
\end{questionbox}

\begin{questionbox}[第 4 题：简单线性回归、最小二乘与残差（18 分）]
对 $n=8$ 个观测拟合含截距的简单线性回归
\[
  Y_i=\beta_0+\beta_1x_i+\varepsilon_i.
\]
已知
\[
  \bar x=5,\qquad
  \bar y=14,\qquad
  S_{xx}=24,\qquad
  S_{xy}=36,\qquad
  S_{yy}=70.
\]

\begin{enumerate}
  \item 求最小二乘估计 $\hat\beta_0,\hat\beta_1$，并写出拟合直线。 \pts{4}
  \item 求 $SSE$、误差方差估计 $s^2$ 和 $R^2$。 \pts{5}
  \item 使用临界值法检验
  \[
    H_0:\beta_1=0
    \qquad\text{vs.}\qquad
    H_a:\beta_1\ne0.
  \]
  \pts{4}
  \item 从最小二乘正规方程证明
  \[
    \sum_{i=1}^n e_i=0,
    \qquad
    \sum_{i=1}^n x_ie_i=0.
  \]
  再说明为什么“原始误差独立”并不意味着“残差独立”。 \pts{5}
\end{enumerate}
\end{questionbox}

\begin{questionbox}[第 5 题：Week 9 回归证明与 MLE 推断（34 分）]
\textbf{第一部分：简单线性回归 Week 9 证明。}

设 $Y_1,\ldots,Y_n$ 相互独立，且
\[
  \E(Y_i)=\beta_0+\beta_1x_i,
  \qquad
  \Var(Y_i)=\sigma^2.
\]
记
\[
  S_{xx}=\sum_{i=1}^n(x_i-\bar x)^2,
  \qquad
  a_i=\frac{x_i-\bar x}{S_{xx}},
  \qquad
  b_i=\frac1n-\bar x a_i.
\]
可直接使用以下恒等式：
\[
  \sum a_i=0,\quad
  \sum a_ix_i=1,\quad
  \sum a_i^2=\frac1{S_{xx}},
\]
\[
  \sum b_i=1,\quad
  \sum b_ix_i=0,\quad
  \sum b_i^2=\frac1n+\frac{\bar x^2}{S_{xx}},
  \quad
  \sum a_ib_i=-\frac{\bar x}{S_{xx}}.
\]

\begin{enumerate}
  \item 已知
  \[
    \hat\beta_1=\sum a_iY_i,
    \qquad
    \hat\beta_0=\sum b_iY_i.
  \]
  证明 $\E(\hat\beta_1)=\beta_1$ 且 $\E(\hat\beta_0)=\beta_0$。 \pts{4}
  \item 证明
  \[
    \Var(\hat\beta_1)=\frac{\sigma^2}{S_{xx}},
    \qquad
    \Var(\hat\beta_0)=\sigma^2\left(\frac1n+\frac{\bar x^2}{S_{xx}}\right),
  \]
  以及
  \[
    \Cov(\hat\beta_0,\hat\beta_1)
    =
    -\frac{\sigma^2\bar x}{S_{xx}}.
  \]
  \pts{5}
  \item 令 $\hat m(x_0)=\hat\beta_0+\hat\beta_1x_0$。证明
  \[
    \E\{\hat m(x_0)\}=\beta_0+\beta_1x_0,
  \]
  且
  \[
    \Var\{\hat m(x_0)\}
    =
    \sigma^2\left[
      \frac1n+\frac{(x_0-\bar x)^2}{S_{xx}}
    \right].
  \]
  \pts{4}
  \item 残差可写成
  \[
    e_i=Y_i-\bar Y-(x_i-\bar x)\hat\beta_1.
  \]
  证明 $\E(e_i)=0$。再利用
  \[
    \Cov(Y_i-\bar Y,\hat\beta_1)
    =
    \frac{(x_i-\bar x)\sigma^2}{S_{xx}}
  \]
  推出
  \[
    \Var(e_i)
    =
    \sigma^2\left[
      1-\frac1n-\frac{(x_i-\bar x)^2}{S_{xx}}
    \right].
  \]
  \pts{5}
\end{enumerate}

\textbf{第二部分：四类分布的 MLE 与推断。}
以下各小题相互独立。

\begin{enumerate}[resume]
  \item 若 $Y_1,\ldots,Y_{10}\iid N(\mu,\sigma^2)$，并给出
  \[
    \bar y=3,
    \qquad
    \sum_{i=1}^{10}(y_i-\bar y)^2=40,
  \]
  求 $\mu$ 与 $\sigma^2$ 的 MLE，并解释为什么 $\hat\sigma^2_{\mathrm{MLE}}$ 与无偏样本方差分母不同。 \pts{4}

  \item 若 $Y\sim\operatorname{Binomial}(80,p)$ 且观测到 $y=20$，求 $\hat p_{\mathrm{MLE}}$ 与其基于 Fisher 信息的大样本标准误。说明 $y=0$ 或 $y=80$ 时应如何检查 MLE。 \pts{4}

  \item 若 $Y_1,\ldots,Y_{25}\iid\operatorname{Poisson}(\lambda)$ 且
  $\sum y_i=100$，求 $\hat\lambda_{\mathrm{MLE}}$、Fisher 信息，并构造 95\% Wald 置信区间。 \pts{4}

  \item 若 $Y_1,\ldots,Y_{20}\iid\operatorname{Exponential}(\lambda)$，其中 $\lambda$ 为 rate，且
  $\sum y_i=50$。求 $\hat\lambda_{\mathrm{MLE}}$，并用 Wald 临界值法检验
  \[
    H_0:\lambda=0.5
    \qquad\text{vs.}\qquad
    H_a:\lambda\ne0.5.
  \]
  \pts{4}
\end{enumerate}
\end{questionbox}

\newpage
% ============================================================
\section{详细解析}
% ============================================================

\begin{checkbox}[阅卷视角下的通用得分点]
\begin{enumerate}
  \item 先定义总体参数和样本设计，再选择标准误与参考分布。
  \item 检验必须写出拒绝域并比较临界值；不能只报告统计量。
  \item Week 9 回归证明题应先写出线性组合形式，再用给定恒等式推出期望、方差和协方差。
  \item MLE 题应写似然或 log-likelihood、score、参数空间或边界检查，再进行 Fisher 信息推断。
\end{enumerate}
\end{checkbox}

\begin{solutionbox}[第 1 题解析：无偏估计与 MSE]
\textbf{（1）证明偏差--方差分解。}

令 $b=\Bias(T)=\E(T)-\theta$。在误差中加减 $\E(T)$：
\[
  T-\theta
  =
  \bigl(T-\E(T)\bigr)
  +
  \bigl(\E(T)-\theta\bigr).
\]
平方并取期望：
\[
\begin{aligned}
  \MSE(T)
  &=
  \E[(T-\theta)^2]\\
  &=
  \E[(T-\E T)^2]
  +2b\,\E[T-\E T]
  +b^2.
\end{aligned}
\]
因为 $\E[T-\E T]=0$，所以交叉项消失：
\[
  \boxed{\MSE(T)=\Var(T)+\Bias(T)^2.}
\]

\textbf{（2）无偏性与最优权重。}

\[
\begin{aligned}
  \E(T_a)
  &=
  a\E(T_1)+(1-a)\E(T_2)\\
  &=
  a\theta+(1-a)\theta
  =
  \theta.
\end{aligned}
\]
因此 $T_a$ 对任意实数 $a$ 都无偏。无偏时 MSE 等于方差。又因为 $T_1,T_2$ 独立，
\[
\begin{aligned}
  \MSE(T_a)
  &=
  \Var(T_a)\\
  &=
  a^2\frac4n+(1-a)^2\frac9n\\
  &=
  \frac{4a^2+9(1-a)^2}{n}.
\end{aligned}
\]
对 $a$ 求导：
\[
  \frac{d}{da}\{4a^2+9(1-a)^2\}
  =
  8a-18(1-a)
  =
  26a-18.
\]
令导数为零：
\[
  \boxed{a^\star=\frac{18}{26}=\frac9{13}.}
\]
二阶导数为 $26>0$，故为最小值。

\textbf{（3）最小 MSE 与解释。}
\[
\begin{aligned}
  \MSE(T_{a^\star})
  &=
  \frac1n\left[
  4\left(\frac9{13}\right)^2
  +
  9\left(\frac4{13}\right)^2
  \right]\\
  &=
  \frac{468}{169n}
  =
  \boxed{\frac{36}{13n}}.
\end{aligned}
\]
而
\[
  \MSE(T_1)=\frac4n,
  \qquad
  \MSE(T_2)=\frac9n.
\]
因为 $36/13\approx2.769<4<9$，最优组合比单独使用任一估计量都更精确。$T_1$ 的方差更小，因此最优组合给它更大的权重 $9/13$。这体现了“精度越高，权重越大”的原则。
\end{solutionbox}

\begin{solutionbox}[第 2 题解析：单总体均值与检验错误]
\textbf{方法选择。}总体正态、总体方差未知，因此使用单样本 $t$ 推断，自由度
\[
  \nu=n-1=15.
\]
标准误为
\[
  \SE(\bar X)=\frac{s}{\sqrt n}
  =
  \frac4{\sqrt{16}}
  =
  1.
\]

\textbf{（1）95\% 置信区间。}
\[
\begin{aligned}
  \bar x
  \pm
  t_{0.025,15}\frac{s}{\sqrt n}
  &=
  52
  \pm
  2.131(1)\\
  &=
  \boxed{(49.869,\ 54.131)}.
\end{aligned}
\]
解释：按此方法重复抽样构造的区间中，约 95\% 会覆盖真实总体均值 $\mu$。

\textbf{（2）单侧检验。}

原假设是复合原假设，但参考分布在边界 $\mu=50$ 处计算：
\[
  t_{\mathrm{obs}}
  =
  \frac{\bar x-50}{s/\sqrt n}
  =
  \frac{52-50}{1}
  =
  2.
\]
在 $H_0$ 边界及正态假设下，
\[
  T\sim t_{15}.
\]
这是右尾检验，拒绝域为
\[
  t_{\mathrm{obs}}>t_{0.05,15}=1.753.
\]
因为
\[
  2>1.753,
\]
所以在 5\% 显著性水平下\textbf{拒绝 $H_0$}。样本提供了足够证据支持总体均值大于 50。

\textbf{（3）两类错误。}
\begin{itemize}
  \item 第一类错误：真实均值实际上不大于 50，却根据样本错误地断言均值大于 50。
  \item 第二类错误：真实均值实际上大于 50，却未能拒绝 $H_0$，没有识别出均值大于 50。
\end{itemize}
\end{solutionbox}

\begin{solutionbox}[第 3 题解析：两总体均值与总体比例]
\textbf{第一部分：两个独立均值。}

两个样本独立、总体正态且假设等方差，因此使用 pooled $t$ 方法。

\textbf{（1）合并方差与标准误。}
\[
\begin{aligned}
  s_p^2
  &=
  \frac{(10-1)6^2+(12-1)5^2}{10+12-2}\\
  &=
  \frac{9(36)+11(25)}{20}
  =
  \boxed{29.95}.
\end{aligned}
\]
\[
  s_p=\sqrt{29.95}=\boxed{5.473}.
\]
\[
\begin{aligned}
  \SE(\bar X_1-\bar X_2)
  &=
  s_p\sqrt{\frac1{10}+\frac1{12}}\\
  &=
  5.473\sqrt{0.18333}
  =
  \boxed{2.343}.
\end{aligned}
\]
自由度为 $10+12-2=20$。

\textbf{（2）95\% 置信区间。}
\[
\begin{aligned}
  (\bar x_1-\bar x_2)
  \pm
  t_{0.025,20}\SE
  &=
  6\pm2.086(2.343)\\
  &=
  \boxed{(1.112,\ 10.888)}.
\end{aligned}
\]
区间整体为正，表示数据支持总体 1 的平均值高于总体 2。

\textbf{（3）双侧检验。}
\[
  t_{\mathrm{obs}}
  =
  \frac{6-0}{2.343}
  =
  2.561.
\]
拒绝域为
\[
  |t_{\mathrm{obs}}|>t_{0.025,20}=2.086.
\]
因为 $2.561>2.086$，所以拒绝 $H_0$；有证据表明两个总体均值不同。

\textbf{第二部分：单总体比例。}

\[
  \hat p=\frac{116}{200}=0.58.
\]

\textbf{（4）95\% Wald 置信区间。}

区间估计使用 $\hat p$ 计算标准误：
\[
  \SE_{\mathrm{CI}}
  =
  \sqrt{\frac{0.58(0.42)}{200}}
  =
  0.03490.
\]
\[
  0.58\pm1.96(0.03490)
  =
  \boxed{(0.512,\ 0.648)}.
\]

\textbf{（5）单侧比例检验。}

检验时在原假设边界 $p_0=0.5$ 下计算标准误：
\[
  \SE_{H_0}
  =
  \sqrt{\frac{0.5(0.5)}{200}}
  =
  0.03536.
\]
\[
  Z_{\mathrm{obs}}
  =
  \frac{0.58-0.5}{0.03536}
  =
  2.263.
\]
右尾拒绝域为
\[
  Z_{\mathrm{obs}}>z_{0.05}=1.645.
\]
因为 $2.263>1.645$，拒绝 $H_0$；数据支持总体支持比例大于 50\%。
\end{solutionbox}

\begin{solutionbox}[第 4 题解析：简单线性回归、最小二乘与残差]
\textbf{（1）最小二乘估计。}
\[
  \hat\beta_1
  =
  \frac{S_{xy}}{S_{xx}}
  =
  \frac{36}{24}
  =
  1.5.
\]
\[
  \hat\beta_0
  =
  \bar y-\hat\beta_1\bar x
  =
  14-1.5(5)
  =
  6.5.
\]
所以拟合直线为
\[
  \boxed{\hat y=6.5+1.5x.}
\]

\textbf{（2）平方和、误差方差与 $R^2$。}
\[
\begin{aligned}
  SSE
  &=
  S_{yy}-\frac{S_{xy}^2}{S_{xx}}\\
  &=
  70-\frac{36^2}{24}\\
  &=
  70-54
  =
  \boxed{16}.
\end{aligned}
\]
\[
  s^2
  =
  \frac{SSE}{n-2}
  =
  \frac{16}{6}
  =
  \boxed{\frac83\approx2.667}.
\]
回归平方和为 $SSR=54$，因此
\[
  R^2
  =
  \frac{SSR}{S_{yy}}
  =
  \frac{54}{70}
  =
  \boxed{0.771}.
\]
在该线性模型与样本范围内，约 77.1\% 的响应总变异由拟合直线解释。

\textbf{（3）斜率检验。}
\[
  \SE(\hat\beta_1)
  =
  \frac{s}{\sqrt{S_{xx}}}
  =
  \sqrt{\frac{s^2}{S_{xx}}}
  =
  \sqrt{\frac{8/3}{24}}
  =
  \frac13.
\]
\[
  t_{\mathrm{obs}}
  =
  \frac{1.5-0}{1/3}
  =
  4.5.
\]
自由度为 $n-2=6$，双侧拒绝域为
\[
  |t_{\mathrm{obs}}|>t_{0.025,6}=2.447.
\]
因为 $4.5>2.447$，拒绝 $H_0$；数据支持总体线性斜率不为零。

\textbf{（4）残差正交性质。}

最小二乘法最小化
\[
  SSE(b_0,b_1)=\sum_{i=1}^n(y_i-b_0-b_1x_i)^2.
\]
记 $e_i=y_i-\hat\beta_0-\hat\beta_1x_i$。分别对 $b_0,b_1$ 求偏导并在最优解处令其为零：
\[
  \frac{\partial SSE}{\partial b_0}
  =
  -2\sum e_i=0
  \quad\Longrightarrow\quad
  \boxed{\sum e_i=0},
\]
\[
  \frac{\partial SSE}{\partial b_1}
  =
  -2\sum x_ie_i=0
  \quad\Longrightarrow\quad
  \boxed{\sum x_ie_i=0}.
\]
矩阵形式下
\[
  \be=(\mathbf I-\mathbf H)\by
  =
  (\mathbf I-\mathbf H)\bvarepsilon,
\]
并且
\[
  \Var(\be)=\sigma^2(\mathbf I-\mathbf H).
\]
其非对角元通常不为零：
\[
  \Cov(e_i,e_j)=-\sigma^2h_{ij}.
\]
所以即使原始误差独立，残差由于共同使用了估计出的回归系数，通常彼此相关，并不独立。
\end{solutionbox}

\begin{solutionbox}[第 5 题解析：Week 9 回归证明与 MLE 推断]
\textbf{第一部分：简单线性回归 Week 9 证明。}

\textbf{（1）估计量的无偏性。}

因为 $\hat\beta_1=\sum a_iY_i$，
\[
\begin{aligned}
  \E(\hat\beta_1)
  &=
  \sum a_i\E(Y_i)\\
  &=
  \sum a_i(\beta_0+\beta_1x_i)\\
  &=
  \beta_0\sum a_i+\beta_1\sum a_ix_i\\
  &=
  0+\beta_1=\boxed{\beta_1}.
\end{aligned}
\]
同理，因为 $\hat\beta_0=\sum b_iY_i$，
\[
\begin{aligned}
  \E(\hat\beta_0)
  &=
  \sum b_i\E(Y_i)\\
  &=
  \beta_0\sum b_i+\beta_1\sum b_ix_i\\
  &=
  \boxed{\beta_0}.
\end{aligned}
\]

\textbf{（2）方差与协方差。}

由于 $Y_i$ 相互独立且 $\Var(Y_i)=\sigma^2$，
\[
\begin{aligned}
  \Var(\hat\beta_1)
  &=
  \Var\left(\sum a_iY_i\right)
  =
  \sum a_i^2\Var(Y_i)\\
  &=
  \sigma^2\sum a_i^2
  =
  \boxed{\frac{\sigma^2}{S_{xx}}}.
\end{aligned}
\]
同理
\[
\begin{aligned}
  \Var(\hat\beta_0)
  &=
  \Var\left(\sum b_iY_i\right)
  =
  \sum b_i^2\Var(Y_i)\\
  &=
  \sigma^2\sum b_i^2
  =
  \boxed{\sigma^2\left(\frac1n+\frac{\bar x^2}{S_{xx}}\right)}.
\end{aligned}
\]
协方差为
\[
\begin{aligned}
  \Cov(\hat\beta_0,\hat\beta_1)
  &=
  \Cov\left(\sum b_iY_i,\sum a_iY_i\right)\\
  &=
  \sum a_ib_i\Var(Y_i)\\
  &=
  \sigma^2\sum a_ib_i
  =
  \boxed{-\frac{\sigma^2\bar x}{S_{xx}}}.
\end{aligned}
\]

\textbf{（3）平均响应估计量。}

由上一问无偏性，
\[
  \E\{\hat m(x_0)\}
  =
  \E(\hat\beta_0)+x_0\E(\hat\beta_1)
  =
  \boxed{\beta_0+\beta_1x_0}.
\]
方差为
\[
\begin{aligned}
  \Var\{\hat m(x_0)\}
  &=
  \Var(\hat\beta_0+x_0\hat\beta_1)\\
  &=
  \Var(\hat\beta_0)+x_0^2\Var(\hat\beta_1)
  +2x_0\Cov(\hat\beta_0,\hat\beta_1)\\
  &=
  \sigma^2\left(\frac1n+\frac{\bar x^2}{S_{xx}}\right)
  +x_0^2\frac{\sigma^2}{S_{xx}}
  -2x_0\frac{\sigma^2\bar x}{S_{xx}}\\
  &=
  \boxed{
  \sigma^2\left[
    \frac1n+\frac{(x_0-\bar x)^2}{S_{xx}}
  \right]}.
\end{aligned}
\]

\textbf{（4）残差期望与方差。}

由题给表示式
\[
  e_i=Y_i-\bar Y-(x_i-\bar x)\hat\beta_1.
\]
先求期望：
\[
\begin{aligned}
  \E(e_i)
  &=
  \E(Y_i)-\E(\bar Y)-(x_i-\bar x)\E(\hat\beta_1)\\
  &=
  (\beta_0+\beta_1x_i)-(\beta_0+\beta_1\bar x)
  -(x_i-\bar x)\beta_1\\
  &=
  \boxed{0}.
\end{aligned}
\]
令 $d_i=x_i-\bar x$。又有
\[
  \Var(Y_i-\bar Y)=\sigma^2\left(1-\frac1n\right),
  \qquad
  \Var(\hat\beta_1)=\frac{\sigma^2}{S_{xx}},
\]
以及题给
\[
  \Cov(Y_i-\bar Y,\hat\beta_1)
  =
  \frac{d_i\sigma^2}{S_{xx}}.
\]
因此
\[
\begin{aligned}
  \Var(e_i)
  &=
  \Var(Y_i-\bar Y-d_i\hat\beta_1)\\
  &=
  \Var(Y_i-\bar Y)+d_i^2\Var(\hat\beta_1)
  -2d_i\Cov(Y_i-\bar Y,\hat\beta_1)\\
  &=
  \sigma^2\left(1-\frac1n\right)
  +d_i^2\frac{\sigma^2}{S_{xx}}
  -2d_i\frac{d_i\sigma^2}{S_{xx}}\\
  &=
  \boxed{
  \sigma^2\left[
    1-\frac1n-\frac{(x_i-\bar x)^2}{S_{xx}}
  \right]}.
\end{aligned}
\]

\textbf{第二部分：四类分布的 MLE 与推断。}

\textbf{（5）正态分布。}

当 $\mu,\sigma^2$ 都未知时，
\[
  \boxed{\hat\mu_{\mathrm{MLE}}=\bar y=3},
\]
\[
  \boxed{
  \hat\sigma^2_{\mathrm{MLE}}
  =
  \frac1n\sum_{i=1}^n(y_i-\bar y)^2
  =
  \frac{40}{10}
  =
  4.
  }
\]
MLE 通过最大化似然得到，因此方差估计分母为 $n$。无偏样本方差为了修正估计均值损失的一个自由度，使用分母 $n-1$：
\[
  S^2=\frac{40}{9}.
\]
所以 $\hat\sigma^2_{\mathrm{MLE}}$ 一般不是 $\sigma^2$ 的无偏估计。

\textbf{（6）二项分布。}

\[
  \boxed{\hat p_{\mathrm{MLE}}=\frac{y}{m}=\frac{20}{80}=0.25.}
\]
单个 $\operatorname{Binomial}(m,p)$ 观测的 Fisher 信息为
\[
  I(p)=\frac{m}{p(1-p)}.
\]
代入 MLE 后，
\[
  \widehat{\SE}(\hat p)
  =
  \sqrt{\frac{\hat p(1-\hat p)}m}
  =
  \sqrt{\frac{0.25(0.75)}{80}}
  =
  \boxed{0.0484}.
\]
若 $y=0$ 或 $y=m$，score 方程在参数空间内部没有根，必须检查 $p=0$ 或 $p=1$ 的边界；结果仍为 $\hat p=y/m$。

\textbf{（7）Poisson 分布。}

Poisson log-likelihood 的 score 为
\[
  U(\lambda)
  =
  \frac{\sum y_i}{\lambda}-n.
\]
令其为零：
\[
  \boxed{
  \hat\lambda
  =
  \bar y
  =
  \frac{100}{25}
  =
  4.
  }
\]
Fisher 信息为
\[
  \boxed{I_n(\lambda)=\frac n\lambda.}
\]
因此
\[
  \widehat{\SE}(\hat\lambda)
  =
  \frac1{\sqrt{I_n(\hat\lambda)}}
  =
  \sqrt{\frac{\hat\lambda}{n}}
  =
  \sqrt{\frac4{25}}
  =
  0.4.
\]
95\% Wald 区间为
\[
  4\pm1.96(0.4)
  =
  \boxed{(3.216,\ 4.784)}.
\]

\textbf{（8）指数分布的 rate 参数。}

指数 rate 参数的 log-likelihood 为
\[
  \ell(\lambda)
  =
  n\log\lambda-\lambda\sum y_i.
\]
score 方程为
\[
  \frac n\lambda-\sum y_i=0,
\]
所以
\[
  \boxed{
  \hat\lambda
  =
  \frac n{\sum y_i}
  =
  \frac{20}{50}
  =
  0.4.
  }
\]
Fisher 信息和估计标准误为
\[
  I_n(\lambda)=\frac n{\lambda^2},
  \qquad
  \widehat{\SE}(\hat\lambda)
  =
  \frac{\hat\lambda}{\sqrt n}
  =
  \frac{0.4}{\sqrt{20}}
  =
  0.0894.
\]
Wald 统计量为
\[
  Z_{\mathrm{obs}}
  =
  \frac{0.4-0.5}{0.0894}
  =
  -1.118.
\]
双侧拒绝域为
\[
  |Z_{\mathrm{obs}}|>z_{0.025}=1.960.
\]
因为 $|-1.118|<1.960$，不能拒绝 $H_0$。现有样本没有提供足够证据表明指数分布的 rate 与 0.5 不同。
\end{solutionbox}

\begin{warnbox}[最后核对：Wald 区间的局限]
第 3 题比例区间和第 5 题第二部分的 MLE 区间均属于 Wald 近似；样本很小、参数接近边界或区间越出参数空间时，需说明近似可能不可靠。
\end{warnbox}

\end{document}
